% !TEX root = one_parameter_lambda.tex
\documentclass[12pt,a4paper]{article}
\usepackage{cmap}
\usepackage[margin=1in]{geometry}
\usepackage{amsmath,amssymb,amsfonts}
\usepackage[utf8]{inputenc}
\usepackage[T2A]{fontenc}
\usepackage[russian]{babel}
\usepackage[hidelinks]{hyperref}
\usepackage{comfortaa}
\renewcommand*\familydefault{\sfdefault}
\usepackage[x11names]{xcolor}

\makeatletter
\renewcommand{\@makefnmark}{\mbox{\textcolor{white}{\@thefnmark}}}
\renewcommand{\@makefntext}[1]{\parindent\fnindent\noindent\mbox{}\llap{\@thefnmark}\textcolor{white}{#1}}
\renewcommand{\footnoterule}{\kern-3pt\color{white}\hrule width 0.4\textwidth height 0.4pt\kern 2.6pt}
\makeatother

\title{\textbf{Однопараметрическая $\Lambda$-космология}}
\author{Олег Понфиленок\thanks{Независимый исследователь; \href{mailto:ponfil@gmail.com}{ponfil@gmail.com}}}
\date{}

\begin{document}

\pagecolor{black}
\color{white}

\maketitle

\begin{abstract}
Стандартная модель $\Lambda$CDM задаёт эволюцию Вселенной набором якобы независимых параметров: $H_0$, $\Omega_{m,0}$, $\Omega_{\Lambda,0}$, $T_0$. Мы показываем, что плоская пылевая Вселенная с космологической постоянной полностью определяется космологической постоянной $\Lambda$ и безразмерным моментом наблюдения $\tau_0$.
Вводя безразмерное время $\tau = \frac{3}{2}H_\Lambda t$ (где $H_\Lambda = c/R_\Lambda$, $R_\Lambda = \sqrt{3/\Lambda}$), вся эволюция и наблюдаемые в эпоху $\tau_0$ записываются через функцию $\sinh(\tau)$: масштабный фактор, параметр Хаббла, плотности компонентов, температура фона, полный причинный радиус.
Вакуумный постулат $\bar\lambda_\Lambda = \delta R_{\max} = \xi\,(R_{\max}\ell_P^2)^{1/3} = \sqrt{R_\Lambda\,\ell_P} = r_s$ в эпоху равенства плотностей $\tau_{m\Lambda}$ однозначно определяет $\xi = \sqrt[3]{2/I}$, где $I = \Gamma(\tfrac{1}{6})\,\Gamma(\tfrac{1}{3})/(3\sqrt{\pi}) \approx 2{,}804$.
При заданной температуре реликтового фона $T_0$ (фиксирующей момент наблюдения $\tau_0$) и доле барионов $\Omega_{b,0}/\Omega_{m,0}$ модель воспроизводит: $H_0 \approx 67$~км/с/Мпк, долю материи $\Omega_{m,0} \approx 0{,}31$ и фотон-барионное отношение $N_\gamma/N_b \approx 1{,}6 \times 10^9$. Кроме того, она даёт голографическую переформулировку проблемы космологической постоянной, выражая наблюдаемую $\rho_\Lambda$ через инфракрасно-ультрафиолетовый масштаб $\bar\lambda_\Lambda = \sqrt{R_\Lambda\,\ell_P}$ вместо планковского, чем устраняет наивное расхождение в $10^{120}$ раз.
Вселенная оказывается столь же простой, как невращающаяся незаряженная чёрная дыра.
\end{abstract}

\noindent\textbf{Ключевые слова:} космологическая постоянная; флуктуации вакуума; неопределённость Каройхази; голографический принцип; фотон-барионное отношение

%%% ===================================================================
\section{Введение: избыточность параметров \texorpdfstring{$\Lambda$CDM}{ΛCDM}}
%%% ===================================================================

Стандартная модель $\Lambda$CDM~\cite{Planck2018} описывает современное состояние Вселенной набором наблюдательных параметров в текущую эпоху: параметр Хаббла $H_0$, доли плотности $\Omega_{m,0}$, $\Omega_{\Lambda,0}$, температура микроволнового фона $T_0$. Эти величины подаются как начальные условия~--- то, что нужно измерить, чтобы зафиксировать конкретную вселенную из семейства возможных решений уравнений Фридмана.

В данной работе мы покажем, что эти параметры избыточны: плоская $\Lambda$CDM-вселенная полностью описывается космологической постоянной $\Lambda$ и безразмерным моментом времени $\tau_0$, в котором мы наблюдаем.

%%% ===================================================================
\section{Обозначения и фундаментальные масштабы}
%%% ===================================================================

Введём ключевые обозначения, определяемые единственным параметром $\Lambda$ и фундаментальными константами $(c, \hbar, G, k_B)$:

\begin{itemize}
\item $t$~--- космологическое время (возраст Вселенной);
\item $R_\Lambda = \sqrt{3/\Lambda}$~--- радиус де Ситтера;
\item $H_\Lambda = c/R_\Lambda$~--- асимптотическая постоянная Хаббла;
\item $\rho_\Lambda = \Lambda c^2/(8\pi G)$~--- плотность тёмной энергии;
\item $T_{\mathrm{dS}} = \hbar H_\Lambda/(2\pi k_B)$~--- температура горизонта де Ситтера;
\item $\ell_P = \sqrt{\hbar G/c^3}$~--- планковская длина;
\item $m_P = \sqrt{\hbar c/G}$~--- планковская масса.
\end{itemize}

Безразмерное время:
\begin{equation}
\boxed{\tau = \frac{3}{2}\,H_\Lambda\,t = \frac{3\,c\,t}{2\,R_\Lambda}.}
\end{equation}

%%% ===================================================================
\section{Эволюция через \texorpdfstring{$\sinh(\tau)$}{sinh(τ)}}
\label{sec:sinh}
%%% ===================================================================

Для плоской вселенной, заполненной пылью и космологической постоянной, точное решение уравнения Фридмана в безразмерном времени $\tau=\tfrac{3}{2}H_\Lambda t$ есть~\cite{Ryden2017} (см.~там же $\S$~5.4.2, ур.~(5.101)). Поскольку после рекомбинации $T \propto a^{-1}$, отсюда вся динамика Вселенной выражается через единственную функцию $\sinh(\tau)$ (коэффициенты $A$ и $B$ будут подставлены в \S~\ref{sec:taumax}):
\begin{align}
a(\tau) &= A\,\sinh^{2/3}(\tau), \label{eq:a_sinh} \\[6pt]
H(\tau) &= H_\Lambda \cdot \coth(\tau), \label{eq:hubble} \\[6pt]
\rho_m(\tau) &= \frac{\rho_\Lambda}{\sinh^2(\tau)}, \label{eq:rho_m} \\[6pt]
T(\tau) &= B\,\sinh^{-2/3}(\tau), \label{eq:temp_sinh} \\[6pt]
\Omega_m(\tau) &= \operatorname{sech}^2(\tau), \label{eq:omega_m} \\[6pt]
\Omega_\Lambda(\tau) &= \tanh^2(\tau). \label{eq:omega_lambda}
\end{align}
Уравнения~(\ref{eq:hubble})--(\ref{eq:omega_lambda}) не содержат свободных параметров: каждая наблюдаемая величина есть функция единственного аргумента $\tau$, определяемого через $\Lambda$ и космологическое время $t$.

%%% ===================================================================
\section{Полный причинный радиус и интеграл \texorpdfstring{$I$}{I}}
%%% ===================================================================

Определим \emph{полный причинный радиус} Вселенной как сумму горизонта частиц и горизонта событий~\cite{Davis2004}:
\begin{equation}
R = R_{\mathrm{p}} + R_{\mathrm{e}} = c\,a\int_0^{\infty}\frac{\mathrm{d}t'}{a(t')}.
\end{equation}
Физический смысл этой величины~--- \emph{максимальный причинно доступный масштаб за всё время наблюдения}. Горизонт частиц $R_{\mathrm{p}}$ задаёт то, что сопутствующий наблюдатель способен увидеть (прошлый световой конус, $\int_0^{t}\mathrm{d}t'/a$), а горизонт событий $R_{\mathrm{e}}$~--- то, с чем он может вступить в причинную связь в будущем (будущий световой конус, $\int_t^{\infty}\mathrm{d}t'/a$). Их сумма~--- это полная сопутствующая область, доступная наблюдателю за всю историю Вселенной, а не мгновенный размер видимой части в фиксированную эпоху. В пылевой Вселенной с $\Lambda$ оба интеграла сходятся, поэтому $R$ конечен и хорошо определён.
Как показано в предыдущем разделе~(\S~\ref{sec:sinh}, формула~(\ref{eq:a_sinh})), $a(\tau)=A\,\sinh^{2/3}(\tau)$. Подставляя $\mathrm{d}t=(2R_\Lambda/3c)\,\mathrm{d}\tau$,
\begin{equation}
\begin{aligned}
R(\tau) &= c\,a(\tau)\int_0^{\infty}\frac{\mathrm{d}t'}{a(t')} \\
&= c\,A\,\sinh^{2/3}(\tau)\int_0^{\infty}\frac{(2R_\Lambda/3c)\,\mathrm{d}\tau'}{A\,\sinh^{2/3}(\tau')} \\
&= \frac{2}{3}\,R_\Lambda\,\sinh^{2/3}(\tau)
   \int_0^{\infty}\frac{\mathrm{d}\tau'}{\sinh^{2/3}(\tau')},
\end{aligned}
\end{equation}
где коэффициент $A$ сократился в числителе и знаменателе интеграла. Отсюда безразмерный интеграл
\begin{equation}
I = \frac{2}{3}\int_0^{\infty}\frac{\mathrm{d}\tau}{\sinh^{2/3}(\tau)}.
\end{equation}
Замена $z=\mathrm{e}^{-\tau}$ ($\tau=-\ln z$, $\mathrm{d}\tau=-\mathrm{d}z/z$) и тождество $\sinh(-\ln z)=(1-z^2)/(2z)$ приводят к бета-интегралу:
\begin{align}
I &= \frac{2}{3}\int_0^1\left[\frac{1-z^2}{2z}\right]^{-2/3}\frac{\mathrm{d}z}{z}
   = \frac{2^{2/3}}{3}\int_0^1 z^{-1/3}(1-z^2)^{-2/3}\,\mathrm{d}z \notag\\
  &= \frac{2^{2/3}}{6}\int_0^1 u^{-2/3}(1-u)^{-2/3}\,\mathrm{d}u
   \quad (u=z^2)
   = \frac{2^{2/3}}{6}\,B\!\left(\tfrac{1}{3},\tfrac{1}{3}\right)
   = \frac{2^{2/3}}{6}\,\frac{\Gamma\!\bigl(\tfrac{1}{3}\bigr)^2}{\Gamma\!\bigl(\tfrac{2}{3}\bigr)}.
\end{align}
Формула удвоения $\Gamma(2z)=2^{2z-1}\Gamma(z)\Gamma(z+\tfrac12)/\sqrt{\pi}$ при $z=\tfrac13$ и соотношение Эйлера $\Gamma(z)\Gamma(1-z)=\pi/\sin(\pi z)$ при $z=\tfrac56$ дают
\begin{equation}
\boxed{I = \frac{\Gamma\!\bigl(\tfrac{1}{6}\bigr)\,\Gamma\!\bigl(\tfrac{1}{3}\bigr)}{3\sqrt{\pi}} \approx 2{,}804.}
\end{equation}
Полный причинный радиус как функция безразмерного времени:
\begin{equation}
\boxed{R(\tau) = R_\Lambda \cdot I \cdot \sinh^{2/3}(\tau).} \label{eq:radius}
\end{equation}
В современную эпоху ($\Omega_{\Lambda,0}/\Omega_{m,0}=\sinh^2\tau_0=2{,}17$, $\tau_0=1{,}18$) из формулы~(\ref{eq:radius}) при данных Planck~2018 получаем
\begin{equation}
\boxed{R_0 = 63{,}7~\text{Gly}}
\end{equation}
(при $I=2{,}804$). С учётом радиационной поправки $I(\Delta)=2{,}766$~--- $R_0=62{,}9$~Gly. Для сравнения, в стандартной плоской $\Lambda$CDM полный причинный радиус $R_{\mathrm{p}}+R_{\mathrm{e}}$ (горизонт частиц плюс горизонт событий~\cite{Davis2004}) при тех же параметрах Planck даёт $R_0^{\mathrm{obs}}=62{,}8\pm 0{,}6$~Gly.

%%% ===================================================================
\section{Учёт радиационной эпохи: поправка к интегралу \texorpdfstring{$I$}{I}}
%%% ===================================================================

В предыдущих разделах полный причинный радиус рассчитывался в рамках пылевой модели (вселенная, заполненная только пылью и космологической постоянной), что является хорошим приближением для поздних эпох. Однако на ранних этапах эволюции Вселенной доминировало излучение (релятивистские частицы: фотоны реликтового фона и нейтрино). Учтем влияние радиационной эпохи на причинный радиус и интеграл $I$.

Закон расширения Фридмана с учётом радиации принимает вид:
\begin{equation}
H^2(a) = H_0^2 \left( \Omega_{r,0}\,a^{-4} + \Omega_{m,0}\,a^{-3} + \Omega_{\Lambda,0} \right),
\end{equation}
где $\Omega_{r,0}$~--- доля излучения в современную эпоху ($a_0=1$). Полный причинный радиус в сопутствующих координатах (при $a_0 = 1$) теперь равен:
\begin{equation}
R = \frac{c}{H_0} \int_0^{\infty} \frac{\mathrm{d}a}{a^2 \sqrt{\Omega_{r,0}\,a^{-4} + \Omega_{m,0}\,a^{-3} + \Omega_{\Lambda,0}}} = \frac{c}{H_0} \int_0^{\infty} \frac{\mathrm{d}a}{\sqrt{\Omega_{r,0} + \Omega_{m,0}\,a + \Omega_{\Lambda,0}\,a^4}}.
\end{equation}
Сделаем стандартную замену переменной $x = a \left(\Omega_{\Lambda,0}/\Omega_{m,0}\right)^{1/3}$. Это даёт:
\begin{equation}
R = R_\Lambda \left(\frac{\Omega_{\Lambda,0}}{\Omega_{m,0}}\right)^{1/3} \cdot I(\Delta),
\end{equation}
где поправленный интеграл $I(\Delta)$ имеет вид:
\begin{equation}
I(\Delta) = \int_0^{\infty} \frac{\mathrm{d}x}{\sqrt{\Delta + x + x^4}},
\end{equation}
а безразмерный параметр $\Delta$, характеризующий влияние радиации, определяется соотношением:
\begin{equation}
\Delta = \frac{\Omega_{r,0}}{\Omega_{m,0}} \left(\frac{\Omega_{\Lambda,0}}{\Omega_{m,0}}\right)^{1/3}.
\end{equation}

Используя данные Planck 2018~\cite{Planck2018} для современной эпохи: $\Omega_{m,0} \approx 0{,}3153$, $\Omega_{\Lambda,0} \approx 0{,}6847$, температура реликтового излучения $T_0 \approx 2{,}7255$~К и эффективное число сортов нейтрино $N_{\mathrm{eff}} \approx 3{,}046$, получаем плотность излучения $\Omega_{r,0} \approx 9{,}21 \times 10^{-5}$. Отсюда находим значение безразмерного параметра:
\begin{equation}
\Delta \approx 3{,}78 \times 10^{-4}.
\end{equation}

Численное интегрирование даёт значение поправленного интеграла:
\begin{equation}
I(\Delta) \approx 2{,}7659,
\end{equation}
что на $1{,}37\%$ меньше базового значения $I \approx 2{,}8044$, полученного без учёта радиации. 

Таким образом, влияние радиационной эпохи приводит к небольшому сжатию полного причинного горизонта Вселенной на ранних этапах на величину около $1{,}4\%$, однако качественная картина однопараметрической космологии и постоянства массы сохраняется.

%%% ===================================================================
\section{Полная масса пыли}
%%% ===================================================================

Покажем, что полная масса пылевой материи во Вселенной постоянна и определяется только $\Lambda$.

Плотность пылевой материи $\rho_m$ в плоской вселенной с космологической постоянной убывает обратно пропорционально объёму: $\rho_m(\tau) = \rho_\Lambda / \sinh^2(\tau)$. Причинный объём Вселенной определяется через причинный радиус $R(\tau)$ из формулы~(\ref{eq:radius}):
\begin{equation}
V(\tau) = \frac{4}{3}\pi\,R^3(\tau) = \frac{4}{3}\pi\,R_\Lambda^3\,I^3\,\sinh^2(\tau).
\end{equation}

Тогда полная масса пылевой материи $M = \rho_m(\tau)\,V(\tau)$ равна:
\begin{equation}
M = \frac{\rho_\Lambda}{\sinh^2(\tau)} \cdot \frac{4}{3}\pi\,R_\Lambda^3\,I^3\,\sinh^2(\tau).
\end{equation}
Подставляя $\rho_\Lambda = \frac{3 c^2}{8\pi G R_\Lambda^2}$, получаем:
\begin{equation}
\boxed{M = \frac{c^2\,R_\Lambda}{2G}\,I^3 = M_\Lambda \cdot I^3,}
\end{equation}
где $M_\Lambda = c^2 R_\Lambda / (2G) \approx 1{,}1 \times 10^{53}$~кг~--- масса Шварцшильда для радиуса де Ситтера $R_\Lambda$, а $I^3 \approx 22{,}05$ --- безразмерная константа. При $R_\Lambda \approx 1{,}65 \times 10^{26}$~м получаем $M \approx 2{,}4 \times 10^{54}$~кг.

Все зависящие от времени сомножители, а также доли плотностей $\Omega_m$ и $\Omega_\Lambda$ сократились. Таким образом, полная масса пылевой материи Вселенной строго постоянна во времени и однозначно определяется космологической постоянной $\Lambda$.


%%% ===================================================================
\section{Вакуумный постулат и голографическое условие}
%%% ===================================================================

Введём гипотезу, связывающую максимальный причинный масштаб с характерной длиной волны флуктуаций вакуума.

\paragraph{Приведённая длина волны.} При сопоставлении с радиусом Шварцшильда $r_s$ естественно использовать \emph{приведённую} (редуцированную) длину волны $\bar\lambda \equiv \lambda/(2\pi)$: условие $2\pi r = \lambda$ записывается как $r = \bar\lambda$. Полная длина волны $\lambda = 2\pi\bar\lambda$.

\paragraph{Неопределённость Каройхази.} Для пространственной области размера $R$ квантовая неопределённость длины (вызванная флуктуациями метрики) в общем виде оценивается как $\delta R \sim (R \cdot \ell_P^2)^{1/3}$. В нашей модели безразмерный коэффициент $\xi$ фиксирует эту оценку точным равенством; для максимального причинного горизонта $R_{\max}$ определим:
\begin{equation}
\delta R_{\max} = \xi \cdot (R_{\max} \cdot \ell_P^2)^{1/3}.
\end{equation}
Коэффициент $\xi$ определяется голографическим условием (см.~ниже); радиус Шварцшильда $r_s = \sqrt{R_\Lambda\,\ell_P}$ входит в постулат без дополнительного множителя.

\paragraph{Вакуумный постулат.} Характерный масштаб квантовых флуктуаций вакуума задаётся приведённой длиной волны $\bar\lambda_\Lambda$. Мы постулируем её точное равенство неопределённости Каройхази для $R_{\max}$ и радиусу Шварцшильда:
\begin{equation}
\boxed{\bar\lambda_\Lambda = \delta R_{\max} = \xi \cdot (R_{\max} \cdot \ell_P^2)^{1/3} = \sqrt{R_\Lambda \cdot \ell_P} = r_s,}
\end{equation}
Связь масштаба $\bar\lambda_\Lambda = \sqrt{R_\Lambda\,\ell_P}$ с инфракрасным размером горизонта восходит к работе Коэна, Каплана и Нельсона~\cite{Cohen1999}, а также обсуждалась Падманабханом~\cite{Padmanabhan2005} в виде соотношения плотностей $\rho_\Lambda \approx \sqrt{\rho_{\mathrm{UV}}\rho_{\mathrm{IR}}}$.

\paragraph{Приведённая длина волны флуктуаций вакуума.}
\begin{equation}
\bar\lambda_\Lambda = \delta R_{\max} = \xi \cdot (R_{\max} \cdot \ell_P^2)^{1/3} = \sqrt{R_\Lambda \cdot \ell_P} = r_s \approx 5{,}2 \times 10^{-5}~\text{м}.
\end{equation}
Полная длина волны $\lambda_\Lambda = 2\pi\bar\lambda_\Lambda \approx 3{,}3 \times 10^{-4}$~м. Приведённая длина $\bar\lambda_\Lambda$ близка к масштабу Планка-Эйнштейна $\ell_{\mathrm{PE}} = (\hbar G / c^3 \Lambda)^{1/4} \approx 37$~мкм (отношение $\bar\lambda_\Lambda/\ell_{\mathrm{PE}} \approx 1{,}4$).
В качестве \emph{гипотетической} иллюстрации отметим, что в измерениях аномального квантового шума в джозефсоновских переходах на основе YBCO~\cite{Hoffmann2004} спад спектральной плотности шума наблюдался при $\nu \approx 0{,}6$--$0{,}8$~ТГц; соответствующая \emph{приведённая} длина волны $\bar\lambda = c/(2\pi\nu) \approx 60$--$80$~мкм оказывается близкой к $\bar\lambda_\Lambda \approx 52$~мкм ($\nu \approx 0{,}92$~ТГц), тогда как $\ell_{\mathrm{PE}} \approx 37$~мкм соответствовал бы завышенной частоте $\approx 1{,}3$~ТГц. Подчеркнём, что интерпретация этого спада как проявления флуктуаций вакуума остаётся спорной и не получила независимого подтверждения, поэтому данное совпадение мы рассматриваем лишь как ориентир, а не как наблюдательное доказательство модели.

\paragraph{Голографическое условие.}
Для вычисления $\xi$ перейдём в особую эпоху $\tau_{m\Lambda}$ равенства плотностей материи и тёмной энергии:
\begin{equation}
\Omega_m(\tau_{m\Lambda}) = \Omega_\Lambda(\tau_{m\Lambda}) \quad\Longrightarrow\quad \sinh(\tau_{m\Lambda}) = 1, \quad \tau_{m\Lambda} = \ln(1 + \sqrt{2}).
\end{equation}

В эту эпоху причинный радиус равен
\begin{equation}
R_{m\Lambda} = R_\Lambda \cdot I \cdot \sinh^{2/3}(\tau_{m\Lambda}) = R_\Lambda \cdot I,
\end{equation}
а температура фона:
\begin{equation}
T_{m\Lambda} = \frac{\hbar c}{2\pi k_B \cdot \xi^3 \cdot I \cdot \sqrt{R_\Lambda \cdot \ell_P}}.
\end{equation}

Наложим требование: температура фона $T_{m\Lambda}$ в эпоху $\tau_{m\Lambda}$ в точности равна температуре Хокинга чёрной дыры с радиусом Шварцшильда $r_s = \bar\lambda_\Lambda = \sqrt{R_\Lambda \cdot \ell_P}$:
\begin{equation}
T_{m\Lambda\mathrm{, Hawking}} = \frac{\hbar c}{4\pi k_B\,r_s} = \frac{\hbar c}{4\pi k_B\,\bar\lambda_\Lambda}.
\end{equation}

Приравнивая $T_{m\Lambda} = T_{m\Lambda\mathrm{, Hawking}}$:
\begin{equation}
\frac{\hbar c}{2\pi k_B \cdot \xi^3 \cdot I \cdot \sqrt{R_\Lambda \cdot \ell_P}} = \frac{\hbar c}{4\pi k_B \cdot \sqrt{R_\Lambda \cdot \ell_P}}.
\end{equation}
Сокращая общие множители, получаем:
\begin{equation}
\xi^3 = \frac{2}{I}.
\end{equation}

Отсюда единственный безразмерный коэффициент модели определяется аналитически:
\begin{equation}
\boxed{\xi = \sqrt[3]{\frac{2}{I}} = \sqrt[3]{\frac{6\sqrt{\pi}}{\Gamma(\tfrac{1}{6})\,\Gamma(\tfrac{1}{3})}} \approx 0{,}894.}
\end{equation}

Идея выражения температуры реликтового фона через геометрическое среднее температур Хокинга и Планка близка к современным подходам Хауга и Татума~\cite{Haug2023}, которые связывают масштабы температур космологического и планковского горизонтов. Однако в нашей работе это отношение выводится из уравнений Фридмана и голографического условия в эпоху равенства плотностей $\tau_{m\Lambda}$.

%%% ===================================================================
\section{Численная оценка и эмпирическое подтверждение}
\label{sec:empirical}
%%% ===================================================================

Подстановка данных миссии Planck 2018~\cite{Planck2018} показывает близкое совпадение температур:
\begin{itemize}
\item Planck+lensing+BAO ($\Lambda \approx 1{,}106 \times 10^{-52}$~м$^{-2}$): $T_{\mathrm{Hawking}} \approx 3{,}532$~К, $T_{m\Lambda} \approx 3{,}553$~К (отклонение \textbf{0{,}59\%}).
\item Planck-only ($\Lambda \approx 1{,}089 \times 10^{-52}$~м$^{-2}$): $T_{\mathrm{Hawking}} \approx 3{,}518$~К, $T_{m\Lambda} \approx 3{,}529$~К (отклонение \textbf{0{,}31\%}).
\end{itemize}

\paragraph{Анализ чувствительности к точности $\Lambda$.}
Измеренное значение космологической постоянной $\Lambda$ имеет относительную экспериментальную погрешность около $\delta\Lambda/\Lambda \approx 1{,}92\%$ (согласно Planck 2018: $\Lambda = (1{,}089 \pm 0{,}021) \times 10^{-52}$~м$^{-2}$). 
При изменении $\Lambda$ меняется отношение плотностей тёмной энергии и материи в современную эпоху $\Omega_{\Lambda,0} / \Omega_{m,0}$. Поскольку физическая плотность материи измерена очень точно ($\Omega_{m,0} H_0^2 = \text{const}$), условие равенства плотностей пыли и вакуума $\rho_m(z_{m\Lambda}) = \rho_\Lambda$ требует, чтобы $(1+z_{m\Lambda})^3 = \rho_\Lambda/\rho_{m,0} \propto \Lambda$. Таким образом, изменение $\Lambda$ напрямую сдвигает красное смещение эпохи равенства $z_{m\Lambda} \propto \Lambda^{1/3}$. Соответственно, температура реликтового излучения в эпоху равенства меняется как $T_{m\Lambda} \propto 1+z_{m\Lambda} \propto \Lambda^{1/3}$.

В то же время температура Хокинга масштабируется как $T_{m\Lambda\mathrm{, Hawking}} \propto R_\Lambda^{-1/2} \propto \Lambda^{1/4}$. Из-за того, что обе температуры возрастают с увеличением $\Lambda$ с близкими степенями, их отношение $\frac{T_{m\Lambda\mathrm{, Hawking}}}{T_{m\Lambda}} \propto \Lambda^{-1/12}$ обладает крайне низкой чувствительностью к вариациям $\Lambda$.

Показатель чувствительности (эластичность) отношения температур к изменению $\Lambda$ равен всего:
\begin{equation}
\frac{\mathrm{d}\ln(T_{m\Lambda\mathrm{, Hawking}}/T_{m\Lambda})}{\mathrm{d}\ln\Lambda} = -\frac{1}{12} \approx -0{,}083.
\end{equation}
Благодаря столь малой чувствительности, при экспериментальной погрешности измерения космологической постоянной $\delta\Lambda/\Lambda \approx 1{,}92\%$, соответствующая неопределённость в определении температур составляет:
\begin{itemize}
\item для температуры Хокинга: $\delta T_{\mathrm{Hawking}} / T_{\mathrm{Hawking}} = \frac{1}{4} \frac{\delta\Lambda}{\Lambda} \approx 0{,}48\%$ (абсолютная погрешность $\approx 0{,}017$~К);
\item для температуры равенства: $\delta T_{m\Lambda} / T_{m\Lambda} = \frac{1}{3} \frac{\delta\Lambda}{\Lambda} \approx 0{,}64\%$ (абсолютная погрешность $\approx 0{,}022$~К);
\item для их отношения: $\delta (T_{\mathrm{Hawking}}/T_{m\Lambda}) / (T_{\mathrm{Hawking}}/T_{m\Lambda}) = \frac{1}{12} \frac{\delta\Lambda}{\Lambda} \approx 0{,}16\%$.
\end{itemize}

Отметим, что наблюдаемое расхождение между температурами составляет всего \textbf{0{,}31\%} (для Planck-only) и \textbf{0{,}59\%} (для комбинированных данных Planck+lensing+BAO). Это расхождение (абсолютная разность $\Delta T \approx 0{,}011$--$0{,}021$~К) существенно меньше, чем экспериментальная погрешность определения самих температур из $\Lambda$ ($\approx 0{,}017$--$0{,}022$~К). Тот факт, что реальное расхождение меньше погрешности измерения, а чувствительность отношения температур к вариациям $\Lambda$ крайне низкая, убедительно доказывает физическую состоятельность, устойчивость и наблюдательную подтверждённость предложенного голографического условия.

Отметим, что в рамках самой предлагаемой однопараметрической космологии отношение температур тождественно равно единице для любой вселенной:
\begin{equation}
\frac{T_{m\Lambda\mathrm{, Hawking}}}{T_{m\Lambda}} = \frac{I\,\xi^3}{2} \equiv 1.
\end{equation}
Оно не зависит от значения $\Lambda$. Зависимость же $\propto \Lambda^{-1/12}$ возникает лишь при внешней верификации через стандартную многопараметрическую $\Lambda$CDM, где космологическая постоянная и плотность материи ($\Omega_{m,0} H_0^2$) рассматриваются как независимые параметры. Малость показателя степени ($-1/12$) гарантирует высокую устойчивость этого равенства к экспериментальным погрешностям.

%%% ===================================================================
\section{Максимальный причинный радиус и эпоха \texorpdfstring{$\tau_{\max}$}{τ\_max}}
\label{sec:taumax}
%%% ===================================================================

Теперь все параметры вычисляются аналитически из космологической постоянной $\Lambda$ и фундаментальных констант.

\paragraph{Нормировка масштабного фактора и температуры.}
Зафиксируем абсолютную нормировку условием термического финиша: в предельно далёкую эпоху $\tau_{\max}$, когда $T(\tau_{\max}) = T_{\mathrm{dS}} \approx 2{,}20 \times 10^{-30}$~К, положим $a(\tau_{\max}) = 1$ (причинный радиус $R = R_{\max}$). В формулах~(\ref{eq:a_sinh}) и~(\ref{eq:temp_sinh}) подставляем
\begin{equation}
A = \sinh^{-2/3}(\tau_{\max}), \qquad B = \frac{T_{\mathrm{dS}}}{A}.
\end{equation}
Тогда масштабный фактор и температура фона записываются как
\begin{align}
a(\tau) &= \left[\frac{\sinh(\tau)}{\sinh(\tau_{\max})}\right]^{2/3}, \label{eq:scale} \\[6pt]
T(\tau) &= \frac{T_{\mathrm{dS}}}{a(\tau)}. \label{eq:temp}
\end{align}

\paragraph{Безразмерное время $\tau_{\max}$.}
Эпоха $\tau_{\max}$ определяется как момент, когда температура реликтового излучения падает до температуры де Ситтера: $T(\tau_{\max}) = T_{\mathrm{dS}}$. Согласно формуле эволюции температуры (\ref{eq:temp}), это соответствует масштабному фактору $a(\tau_{\max}) = 1$. Тогда из закона $T(\tau) \propto 1/a(\tau)$ следует, что:
\begin{equation}
T_{\mathrm{dS}} = T_{m\Lambda} \cdot a(\tau_{m\Lambda}),
\end{equation}
где $\tau_{m\Lambda}$~--- эпоха равенства плотностей, в которую $\sinh(\tau_{m\Lambda}) = 1$, а значит $a(\tau_{m\Lambda}) = \sinh^{-2/3}(\tau_{\max})$. Подставляя значение $T_{m\Lambda}$ из подраздела голографического условия и выражая $\sinh^{2/3}(\tau_{\max})$, получаем:
\begin{equation}
\sinh^{2/3}(\tau_{\max}) = \frac{T_{m\Lambda}}{T_{\mathrm{dS}}} = \frac{1}{\xi^3\,I}\,\sqrt{\frac{R_\Lambda}{\ell_P}}.
\end{equation}
С учётом голографического условия $\xi^3 = 2/I$ окончательно имеем:
\begin{equation}
\boxed{\sinh^{2/3}(\tau_{\max}) = \frac{1}{2}\,\sqrt{\frac{R_\Lambda}{\ell_P}} \approx 1{,}6 \times 10^{30},}
\end{equation}
что даёт $\tau_{\max} \approx 105$. Это соответствует возрасту Вселенной:
\begin{equation}
t_{\max} = \frac{2\,R_\Lambda\,\tau_{\max}}{3\,c} \approx 1{,}2 \times 10^{12}~\text{лет}
\end{equation}
--- термический финиш Вселенной, когда реликтовый фон охладится до температуры горизонта де Ситтера. Это примерно в 90 раз больше текущего возраста Вселенной.

\paragraph{Максимальный причинный радиус.}
Соответствующий максимальный физический причинный радиус Вселенной в момент $\tau_{\max}$ равен:
\begin{equation}
\boxed{R_{\max} = R(\tau_{\max}) = R_\Lambda \cdot I \cdot \sinh^{2/3}(\tau_{\max}) = \frac{I}{2}\,\sqrt{\frac{R_\Lambda^3}{\ell_P}} \approx 7{,}5 \times 10^{56}~\text{м}.}
\end{equation}

%%% ===================================================================
\section{Наблюдаемые параметры нашей эпохи}
%%% ===================================================================

Модель даёт точные значения космологических параметров для нашей эпохи. Космологическая постоянная $\Lambda$ полностью задаёт кривую эволюции температуры $T(\tau)$ (через нормировку $a(\tau_{\max})=1$ при $T(\tau_{\max})=T_{\mathrm{dS}}$, см.~\S~\ref{sec:taumax}). Момент наблюдения $\tau_0$ фиксируется единственным измеренным сегодня масштабом~--- температурой реликтового фона $T_0 = 2{,}7255$~К~\cite{Planck2018}. Из закона эволюции температуры~(\ref{eq:temp}), $T(\tau_0) = T_{\mathrm{dS}}\,[\sinh(\tau_{\max})/\sinh(\tau_0)]^{2/3} = T_0$, откуда
\begin{equation}
\sinh(\tau_0) = \sinh(\tau_{\max})\left(\frac{T_{\mathrm{dS}}}{T_0}\right)^{3/2}, \qquad \tau_0 = 1{,}18.
\end{equation}
Подставляя в формулы эволюции:

\begin{center}
\begin{tabular}{l c c}
\hline
Параметр & Модель & Наблюдения \\
\hline
$\tau_0$ & $1{,}18$ & --- \\
$t_0$ & $13{,}8$~млрд лет & $13{,}8 \pm 0{,}02$~млрд лет \\
$H_0$ & $67$~км/с/Мпк & $67{,}4 \pm 0{,}5$~км/с/Мпк \\
$R_0$ & $63{,}7$~Gly & $62{,}8 \pm 0{,}6$~Gly \\
$T_0$ (вход) & $2{,}7255$~К & $2{,}7255 \pm 0{,}0006$~К \\
$\Omega_{m,0}$ & $0{,}31$ & $0{,}315 \pm 0{,}007$ \\
$\Omega_{\Lambda,0}$ & $0{,}69$ & $0{,}685 \pm 0{,}007$ \\
\hline
\end{tabular}
\end{center}

Единственным входом нашей эпохи служит температура $T_0$, фиксирующая $\tau_0$; все остальные величины в таблице~--- предсказания. В частности, доля материи $\Omega_{m,0} = \operatorname{sech}^2(\tau_0) \approx 0{,}31$ не подгоняется, а следует из $\tau_0$ и согласуется с наблюдаемым значением $0{,}315 \pm 0{,}007$. Сама же кривая $T(\tau)$ задана только $\Lambda$ (нормировка $a(\tau_{\max})=1$ при $T(\tau_{\max})=T_{\mathrm{dS}}$), так что выбор $T_0$ означает лишь указание точки на этой кривой~--- момента, в который мы наблюдаем.

%%% ===================================================================
\section{Сохранение числа частиц в причинном объёме}
%%% ===================================================================

Полное число фотонов и барионов в причинном объёме $V(\tau) = \frac{4\pi}{3}\,R(\tau)^3$ постоянно во времени. Действительно:
\begin{itemize}
\item плотность пыли падает как $\rho_m \propto \sinh^{-2}(\tau)$;
\item плотность фотонов $n_\gamma \propto T^3 \propto a^{-3} \propto \sinh^{-2}(\tau)$;
\item причинный объём растёт как $V \propto R^3 \propto \sinh^2(\tau)$.
\end{itemize}
Все зависимости от $\tau$ взаимно сокращаются, и абсолютные числа частиц в причинном объёме суть чистые числа, определяемые только $R_\Lambda / \ell_P$.

\subsection{Число фотонов}
\begin{equation}
\boxed{N_\gamma = \frac{\zeta(3)\,I^3}{24\pi^4} \left(\frac{R_\Lambda}{\ell_P}\right)^{3/2} \approx 3{,}73 \times 10^{89}.} \label{eq:n_gamma}
\end{equation}

\subsection{Число барионов}
\begin{equation}
\boxed{N_b = \frac{\Omega_b}{\Omega_m}\,\frac{I^3}{2}\,\frac{m_P}{m_p}\,\frac{R_\Lambda}{\ell_P} \approx 2{,}30 \times 10^{80},} \label{eq:n_b}
\end{equation}
при $\Omega_{b,0}/\Omega_{m,0} \approx 0{,}156$ (из данных Planck 2018~\cite{Planck2018}) и $m_b \approx m_p$~--- массе протона.

\subsection{Фотон-барионное отношение}
\begin{equation}
\boxed{\frac{N_\gamma}{N_b} = \frac{\zeta(3)}{12\pi^4}\,\frac{\Omega_m}{\Omega_b}\,\frac{m_p}{m_P}\,\sqrt{\frac{R_\Lambda}{\ell_P}} \approx 1{,}62 \times 10^{9}.} \label{eq:ratio_particles}
\end{equation}
Это отношение прекрасно согласуется с наблюдаемым значением барионной асимметрии Вселенной $\eta_{\mathrm{obs}}^{-1} \approx 1{,}63 \times 10^9$ (соответствующим $\eta \approx 6{,}12 \times 10^{-10}$~\cite{Planck2018}) с точностью лучше $1\%$.

%%% ===================================================================
\section{Голографическая переформулировка проблемы космологической постоянной}
%%% ===================================================================

Плотность тёмной энергии $\rho_\Lambda$ может быть выражена через приведённую длину волны флуктуаций $\bar\lambda_\Lambda$:
\begin{equation}
\rho_\Lambda = C \frac{\hbar}{c\,\bar\lambda_\Lambda^4}.
\end{equation}
Выведем аналитическое выражение для коэффициента $C$. Плотность тёмной энергии связана с космологической постоянной соотношением $\rho_\Lambda = \frac{\Lambda c^2}{8\pi G}$. Выражая $\Lambda$ через радиус де Ситтера $R_\Lambda = \sqrt{3/\Lambda}$, имеем:
\begin{equation}
\rho_\Lambda = \frac{3 c^2}{8\pi G R_\Lambda^2}.
\end{equation}
Используя $\bar\lambda_\Lambda = \sqrt{R_\Lambda \ell_P}$ и $\ell_P^2 = \frac{\hbar G}{c^3}$, запишем:
\begin{equation}
\bar\lambda_\Lambda^4 = R_\Lambda^2 \ell_P^2 = R_\Lambda^2 \frac{\hbar G}{c^3}.
\end{equation}
Выражая отсюда $G R_\Lambda^2 = \bar\lambda_\Lambda^4 c^3/\hbar$ и подставляя в формулу для $\rho_\Lambda$, получаем:
\begin{equation}
\rho_\Lambda = \frac{3}{8\pi} \frac{\hbar}{c\,\bar\lambda_\Lambda^4}.
\end{equation}
Таким образом, точный аналитический коэффициент равен:
\begin{equation}
\boxed{C = \frac{3}{8\pi} \approx 0{,}1194.}
\end{equation}
Следовательно, плотность вакуума выражается через приведённую длину волны $\bar\lambda_\Lambda$ следующим образом:
\begin{equation}
\boxed{\rho_\Lambda = \frac{3}{8\pi} \frac{\hbar}{c\,\bar\lambda_\Lambda^4}.} \label{eq:rho_lambda_final}
\end{equation}
Если перейти к плотности энергии вакуума $E_\Lambda = \rho_\Lambda c^2$, то получим точное квантовое соотношение:
\begin{equation}
\boxed{E_\Lambda = \frac{3}{8\pi} \frac{\hbar c}{\bar\lambda_\Lambda^4}.}
\end{equation}
Численно, при $\bar\lambda_\Lambda \approx 52$~мкм (из данных Planck), эти формулы воспроизводят наблюдаемую $\rho_\Lambda$.

Это соотношение \emph{переформулирует} проблему космологической постоянной в голографическом ключе: наблюдаемая $\rho_\Lambda$ естественно выражается через инфракрасно-ультрафиолетовый масштаб $\bar\lambda_\Lambda = \sqrt{R_\Lambda\,\ell_P}$, и наивное расхождение в $10^{120}$ раз, возникающее при обрезании на планковском масштабе $\ell_P$, не появляется. Подчеркнём, что $\Lambda$ остаётся единственным входным параметром (формула~(\ref{eq:rho_lambda_final}) есть точное переписывание $\rho_\Lambda = 3c^2/(8\pi G R_\Lambda^2)$), однако масштаб обрезания $\bar\lambda_\Lambda$ \emph{не произволен}. Его однозначно задаёт неопределённость Каройхази: $\bar\lambda_\Lambda = \delta R_{\max} = \xi\,(R_{\max}\ell_P^2)^{1/3}$~--- это минимальная геометрически разрешимая длина для максимального причинного горизонта $R_{\max}$. Вакуум обрезается не на планковской длине $\ell_P$, а на этом минимальном масштабе геометрии, допустимом для области размера $R_{\max}$,~--- что согласуется с подходом Коэна--Каплана--Нельсона~\cite{Cohen1999} и Падманабхана~\cite{Padmanabhan2005} и именно поэтому воспроизводит наблюдаемую $\rho_\Lambda$.

Отметим, что в современной теоретической физике отсутствует единое общепринятое значение для безразмерного коэффициента связи между плотностью энергии вакуума и характерной длиной волны квантового обрезания $\lambda$. Обычно в расчетах квантовой теории поля это соотношение имеет вид $E_{\text{vac}} \sim \hbar c / \lambda^4$, где коэффициент $C$ зависит от выбранной схемы регуляризации. Например, для жесткого импульсного обрезания свободного электромагнитного поля на масштабе $k_{\text{max}} = \pi/\lambda$ (что соответствует пределу Найквиста для решеточного шага $\lambda/2$) коэффициент составляет $C = \pi^2/4 \approx 2{,}47$, в то время как при $k_{\text{max}} = \pi/(2\lambda)$ он равен $C = \pi^2/64 \approx 0{,}154$, а для скалярного поля при $k_{\text{max}} = 1/\lambda$ составляет $C = 1/(8\pi^2) \approx 0{,}013$. В геометрически ограниченных системах, таких как металлический куб со стороной $L$ в эффекте Казимира, плотность энергии вакуума определяется точным коэффициентом $C \approx 0{,}092$. Наш аналитически выведенный коэффициент $C = 3/(8\pi) \approx 0{,}119$ естественным образом укладывается в физический диапазон КТП и эффекта Казимира ($10^{-2} \div 10^{-1}$), при этом он не является подгоночным. В отличие от размерного масштаба Бека и де Матоса~\cite{Beck2007}, где коэффициент принимается равным единице по определению (что дает масштаб $37$~мкм и частоту среза шума $\approx 1{,}3$~ТГц), наше значение $\bar\lambda_\Lambda \approx 52$~мкм соответствует более низкой частоте $\nu \approx 0{,}92$~ТГц. Если спад аномального квантового шума в YBCO-переходах~\cite{Hoffmann2004} действительно связан с флуктуациями вакуума~--- что остаётся непроверенной и дискуссионной гипотезой,~--- то наш масштаб согласуется с ним лучше; однако мы не опираемся на этот эксперимент как на подтверждение модели.

%%% ===================================================================
\section{Заключение}
%%% ===================================================================

Под вакуумным постулатом $\bar\lambda_\Lambda = \delta R_{\max} = \xi\,(R_{\max}\ell_P^2)^{1/3} = \sqrt{R_\Lambda\,\ell_P} = r_s$ и голографическим условием в эпоху равенства плотностей ($\xi = \sqrt[3]{2/I}$) плоская пылевая $\Lambda$CDM-вселенная оказывается строго однопараметрической по динамике. Ниже разобраны следствия для наблюдаемых величин и для классических космологических проблем.

\subsection{Однопараметричность}

Единственный фундаментальный параметр $\Lambda$ задаёт массу пыли, $H(\tau)$, доли $\Omega_m(\tau)$ и $\Omega_\Lambda(\tau)$, причинный радиус $R(\tau)$, температуру реликтового фона $T(\tau)$, числа $N_\gamma$ и $N_b$, а также $\tau_{\max} \approx 105$. Четыре привычных «входа» $\Lambda$CDM~--- $H_0$, $\Omega_{m,0}$, $\Omega_{\Lambda,0}$, $T_0$~--- не независимые начальные условия, а значения тех же функций в точке $\tau = \tau_0$. Вместо четырёх свободных параметров остаются две величины: $\Lambda$ и момент наблюдения $\tau_0$.

\subsection{Проблема космологической постоянной}

Стандартное расхождение в $10^{120}$ раз между наивным вакуумным расчётом (с планковским масштабом флуктуаций) и наблюдаемой плотностью $\rho_\Lambda$ в нашей конструкции \emph{переформулируется}, а не разрешается из первых принципов: если характерный масштаб флуктуаций вакуума принять равным не $\ell_P$, а $\bar\lambda_\Lambda = \sqrt{R_\Lambda\,\ell_P} = r_s \approx 52$~мкм, то наблюдаемая $\rho_\Lambda$ воспроизводится автоматически. Это голографическая (ИК--УФ) переформулировка в духе~\cite{Cohen1999,Padmanabhan2005}, причём масштаб обрезания не выбирается произвольно: его задаёт неопределённость Каройхази $\bar\lambda_\Lambda = \delta R_{\max}$~--- минимальная разрешимая геометрия максимального причинного горизонта. Поэтому обрезание вакуума на этом, а не на планковском, масштабе есть следствие квантовой неопределённости пространственной длины, а не подгонка.

\subsection{Фотон-барионное отношение}

Из постоянства полных чисел фотонов и барионов в причинном объёме (формулы~(\ref{eq:n_gamma})--(\ref{eq:n_b})) следует фотон-барионное отношение~(\ref{eq:ratio_particles}). Оно выводится аналитически из $\Lambda$ через $R_\Lambda = \sqrt{3/\Lambda}$ и не является независимым параметром модели: при фиксированной $\Lambda$ и измеренной доле барионов $\Omega_{b,0}/\Omega_{m,0}$ (Planck~2018) получаем $N_\gamma/N_b \approx 1{,}62 \times 10^9$, в согласии с наблюдаемой барионной асимметрией $\eta_{\mathrm{obs}}^{-1} \approx 1{,}63 \times 10^9$ ($\eta \approx 6{,}1 \times 10^{-10}$) с точностью лучше $1\%$.

\subsection{Проблема тонкой настройки}

Проблема отсутствует: при фиксированной $\Lambda$ нет произвольного выбора $H_0$, $\Omega_{m,0}$ или $\Omega_{\Lambda,0}$~--- они однозначно определяются $\tau_0$ (и наоборот, по измеренному $\Omega_{m,0}$ находится $\tau_0$). Аналогично фиксируются $T_0$, $N_\gamma$ и $N_b$.

\subsection{Аналогия с чёрной дырой}

Вселенная по динамике столь же проста, как невращающаяся чёрная дыра Шварцшильда, задаваемая одной массой $M$: здесь эту роль играет $\Lambda$. Эпоха наблюдения $\tau_0$ играет роль «где на орбите смотрим»~--- не новой физики, а выбора момента на уже заданном решении.

\begin{equation}
\boxed{\text{Вселенная} \equiv \Lambda.}
\end{equation}


\bibliographystyle{unsrt}
\bibliography{refs}

\end{document}
